% Main document. Keep this file to preamble + \input of sections so the
% venue template can be swapped by editing only this file.
% Current target: Brazilian law journal / CONPEDI (ABNT layout via abnTeX2).
% See templates/README.md for the SBC (WICS/CSBC) alternative.
\documentclass[article,12pt,oneside,a4paper,brazilian]{abntex2}

% --- Encoding, fonts ------------------------------------------------------
\usepackage[utf8]{inputenc}
\usepackage[T1]{fontenc}
\usepackage{lmodern}
\usepackage{microtype}

% --- Tables, floats -------------------------------------------------------
\usepackage{booktabs}
\usepackage{tabularx}
\usepackage{array}
\usepackage{float}

% --- Bibliography: ABNT author-date via biblatex-abnt ---------------------
% abnTeX2 + biblatex-abnt are compatible; do NOT load abntex2cite.
\usepackage[autostyle]{csquotes}
\usepackage[
  style=abnt,
  backend=biber,
  uniquename=init,
  giveninits,
  justify,
  indent,
  noslsn,
]{biblatex}
\addbibresource{references.bib}

% --- Links ----------------------------------------------------------------
\usepackage{xcolor}
\hypersetup{
  colorlinks=true,
  linkcolor=black,
  citecolor=blue!60!black,
  urlcolor=blue!60!black,
}

% --- Draft helpers --------------------------------------------------------
% \todo{...} prints a visible red marker; redefine as \@gobble before submission.
\newcommand{\todo}[1]{\textcolor{red}{\textbf{[TODO: #1]}}}

% --- Metadata (abnTeX2 commands) ------------------------------------------
\titulo{Da eliminação à indistinguibilidade: um padrão técnico-jurídico para o
  desaprendizado de máquina sob a LGPD e o RGPD}
% abnTeX2 article mode prints the foreign-language title under the main one.
\tituloestrangeiro{From erasure to indistinguishability: a technical-legal
  standard for machine unlearning under the LGPD and the GDPR}
\autor{\todo{autor(es), afiliação, e-mail}}
\local{Brasil}
\data{\today}

\begin{document}

\selectlanguage{brazilian}
\frenchspacing

\maketitle

\begin{resumo}
\todo{Resumo em português (150--250 palavras): problema, tese, método, contribuição.}

\vspace{\onelineskip}
\noindent
\textbf{Palavras-chave}: desaprendizado de máquina; direito à eliminação; LGPD; anonimização; RGPD.
\end{resumo}

\begin{resumo}[Abstract]
\begin{otherlanguage*}{english}
\todo{Abstract in English.}

\vspace{\onelineskip}
\noindent
\textbf{Keywords}: machine unlearning; right to erasure; LGPD; anonymization; GDPR.
\end{otherlanguage*}
\end{resumo}

\textual

\section{Introdução}

% Objetivo da seção: apresentar o problema (modelos treinados retêm informação
% sobre titulares; "eliminar" não tem sentido técnico óbvio), enunciar a tese
% da indistinguibilidade e antecipar a estrutura do artigo.

Esta seção deve situar o leitor no descompasso entre o direito à eliminação,
pensado para bancos de dados, e os modelos de aprendizado de máquina, que não
armazenam registros mas parâmetros influenciados por eles. Deve enunciar a tese
central --- eliminar dados de um modelo significa tornar a contribuição do
titular indistinguível da de um modelo que nunca o viu, julgada contra o
atacante que cada lei pressupõe --- e justificar a centralidade da LGPD: os
arts.~16 e 18, IV oferecem encaixe natural para o padrão, e o direito europeu
(RGPD, EDPB, TJUE) entra como contraponto que confirma a relatividade do
conceito de dado pessoal.

\todo{Gancho inicial: caso concreto ou decisão (EDPS v SRB? Parecer 28/2024?).}
\todo{Delimitar escopo: modelos discriminativos e generativos; excluir bancos de dados vetoriais?}
\todo{Enunciar contribuições em lista (3--4 itens).}
\todo{Roteiro das seções.}

\section{Fundamentos técnicos do desaprendizado de máquina}

% Objetivo da seção: dar ao leitor jurídico o vocabulário mínimo e ao leitor
% técnico o mapa das famílias de métodos e das garantias que cada uma oferece.

Esta seção deve explicar o que significa um modelo "conter" dados de treino
(memorização, inferência de pertinência, extração), apresentar a definição
formal de desaprendizado exato e aproximado, e percorrer as famílias de
métodos: retreinamento completo, SISA \parencite{bourtoule2021machine},
remoção certificada, heurísticas aproximadas e supressão de saídas. Deve
enfatizar que a garantia de cada família é relativa a um \emph{avaliador}
(auditor, atacante) e que avaliações frouxas produzem falsa sensação de
privacidade \parencite{hayes2024inexact}, além dos limites conceituais
apontados por \textcite{cooper2024unlearning} para modelos generativos.

\todo{Definição formal de $(\varepsilon,\delta)$-desaprendizado e relação com privacidade diferencial.}
\todo{Ataques de inferência de pertinência como instrumento de verificação.}
\todo{Benchmarks (TOFU) e por que métricas de utilidade não bastam.}
\todo{Tabela-resumo das famílias: custo, garantia, verificabilidade.}

\section{Marco jurídico: a LGPD e o contraponto europeu}

% Objetivo da seção: reconstruir o direito à eliminação na LGPD como eixo
% principal, identificar o "atacante" que a lei pressupõe ao definir dado
% anonimizado, e só então usar o direito europeu como contraponto.

Esta seção deve partir da LGPD. Primeiro, o art.~12 e o teste de "meios
técnicos razoáveis e disponíveis na ocasião do tratamento", que fixa o
adversário contra o qual a anonimização é medida. Depois, o art.~18 e seu
inciso~IV, que coloca anonimização, bloqueio e eliminação como remédios
alternativos à disposição do titular. Por fim, o art.~16, que autoriza a
conservação de dados após o término do tratamento quando a eliminação for
tecnicamente inviável, cláusula sem equivalente direto no RGPD. Deve incluir a
posição da ANPD (guia de anonimização; eventual regulamentação dos
arts.~18 a 20) e a doutrina brasileira sobre dado anonimizado e sobre os
limites do art.~16.

Em seguida, como contraponto, o direito europeu: art.~17 do RGPD,
Considerando~26 (critério dos "meios razoavelmente suscetíveis de serem
utilizados"), Parecer~28/2024 do EDPB sobre quando um modelo de IA pode ser
considerado anônimo, e o acórdão do TJUE em EDPS~v~SRB (C-413/23~P) sobre a
relatividade do conceito de dado pessoal. O objetivo do contraponto é
mostrar que o adversário relativo já é reconhecido na Europa e que a LGPD,
ao prever o art.~16 e o art.~18, IV, está em posição melhor para acolher o
padrão proposto.

\todo{Texto literal dos arts.~12, 16 e 18 em notas ou apêndice.}
\todo{ANPD: guia de anonimização e regulamentação dos arts.~18--20 (ver notes/perguntas-em-aberto.md).}
\todo{Doutrina brasileira: ver lista de candidatos em notes/perguntas-em-aberto.md.}
\todo{Doutrina europeia: Juliussen et al.; Colonna e Oechtering.}

\section{Quadro comparativo: métodos técnicos e remédios jurídicos}

% Objetivo da seção: cruzar as famílias de métodos da Seção 2 com os remédios
% jurídicos da Seção 3. A tabela é GERADA a partir de notes/crosswalk.csv via
% scripts/crosswalk_to_latex.py; não editar paper/tables/crosswalk.tex à mão.

Esta seção deve apresentar e comentar o quadro comparativo
(Tabela~\ref{tab:crosswalk}), mostrando para cada família de método qual
remédio jurídico ela é capaz de satisfazer, sob quais condições e com qual
grau de verificabilidade. Deve destacar as assimetrias entre RGPD e LGPD
(por exemplo, a ausência no RGPD de uma cláusula equivalente ao art.~16 da
LGPD) e preparar o terreno para a proposta da Seção~5.

% GENERATED FILE - do not edit by hand.
% Source: notes/crosswalk.csv  Generator: scripts/crosswalk_to_latex.py
\begin{table}[htbp]
  \centering
  \small
  \caption{Quadro comparativo entre famílias de métodos de desaprendizado e remédios jurídicos}
  \label{tab:crosswalk}
  \begin{tabularx}{\textwidth}{>{\raggedright\arraybackslash}p{3cm}XXXXX}
    \toprule
    \textbf{Família de método} & \textbf{Eliminação} & \textbf{Anonimização} & \textbf{Bloqueio} & \textbf{Limites técnicos (art. 16)} & \textbf{Esforço desproporcional} \\
    \midrule
    \textbf{Retreinamento completo} &  &  &  &  &  \\
    \textbf{SISA (Sharded Isolated Sliced Aggregated)} &  &  &  &  &  \\
    \textbf{Remoção certificada} &  &  &  &  &  \\
    \textbf{Métodos heurísticos aproximados} &  &  &  &  &  \\
    \textbf{Supressão de saídas (output filtering)} &  &  &  &  &  \\
    \bottomrule
  \end{tabularx}
\end{table}


\todo{Comentário célula a célula após preenchimento do CSV.}
\todo{Critérios usados para classificar cada célula (legenda).}

\section{Proposta: o padrão da indistinguibilidade sob a LGPD}

% Objetivo da seção: formular o padrão técnico-jurídico em linguagem da LGPD
% e mostrar que ele decorre dos arts. 12, 16 e 18, IV; o RGPD entra só para
% demonstrar compatibilidade, não como fundamento.

Esta seção deve enunciar o padrão: a eliminação de dados pessoais de um
modelo está satisfeita quando a contribuição do titular se torna
indistinguível da de um modelo que nunca viu esses dados, avaliada contra o
adversário que o art.~12 da LGPD pressupõe (meios técnicos razoáveis e
disponíveis na ocasião). Deve mostrar que esse resultado é, juridicamente,
uma anonimização, e portanto um dos remédios que o art.~18, IV já coloca ao
lado da eliminação; e que o art.~16 fornece a válvula de escape quando nem
a indistinguibilidade for tecnicamente viável, mediante conservação
documentada e com finalidade limitada. Deve especificar os componentes
operacionais: modelo de adversário, métrica de verificação, ônus da prova
do controlador e documentação exigível pela ANPD. Ao final, deve indicar
brevemente que o mesmo padrão é compatível com o Considerando~26 do RGPD e
com EDPS~v~SRB, embora o RGPD careça de cláusula equivalente ao art.~16.

\todo{Formulação em linguagem normativa (proposta de enunciado interpretativo do art.~18, IV).}
\todo{Como escolher o adversário à luz do art.~12: capacidades, acesso ao modelo, dados auxiliares.}
\todo{Relação com os três testes do Parecer 28/2024 como parâmetro de adversário.}
\todo{Ônus da prova (art.~6, X, responsabilização e prestação de contas) e papel de auditoria independente.}

\section{Objeções e respostas}

% Objetivo da seção: antecipar as críticas mais fortes ao padrão e responder a
% cada uma, sem enfraquecer a tese.

Esta seção deve enfrentar ao menos as seguintes objeções: (i) o padrão
rebaixa a eliminação a uma mera probabilidade, contrariando o texto legal;
(ii) o adversário é indeterminado e muda com o tempo, tornando o padrão
instável; (iii) modelos generativos reproduzem conteúdo sem que isso seja
"dado pessoal" no sentido técnico \parencite{cooper2024unlearning}; (iv) o
padrão favorece grandes controladores com capacidade de auditoria; (v) a
anonimização do art.~18, IV não se aplica a modelos, apenas a bases.

\todo{Listar objeções adicionais surgidas nos fichamentos.}
\todo{Para cada objeção: formulação forte, resposta, concessão residual.}

\section{Recomendações}

% Objetivo da seção: traduzir o padrão em recomendações para os atores
% brasileiros: ANPD, controladores sujeitos à LGPD e legislador (PL 2338/2023).

Esta seção deve propor recomendações concretas e numeradas, cada uma com
destinatário e fundamento: para a ANPD, ao regulamentar os arts.~18 a 20 e
ao consolidar o guia de anonimização, adotar a indistinguibilidade como
critério de atendimento ao pedido de eliminação em modelos; para
controladores sujeitos à LGPD, manter documentação de treino, arquitetura
preparada para desaprendizado e protocolos de verificação compatíveis com
o art.~12; para o legislador, no marco legal da IA (PL 2338/2023), remeter
ao padrão em vez de prescrever técnicas específicas.

\todo{Recomendações numeradas, cada uma com destinatário e fundamento.}
\todo{Verificar estado do PL 2338/2023 antes de redigir (notes/perguntas-em-aberto.md).}

\section{Conclusão}

% Objetivo da seção: recapitular a tese, as contribuições e os limites, e
% apontar a continuação empírica (Proposta A, pasta experiments/).

Esta seção deve retomar a tese em uma frase, resumir o percurso argumentativo,
reconhecer os limites do trabalho (ausência de validação empírica, foco em dois
ordenamentos, dependência de literatura técnica em rápida evolução) e indicar
os próximos passos, em especial a continuação empírica que avaliará métodos
de desaprendizado contra adversários concretos.

\todo{Uma frase de fechamento que conecte técnica e direito.}


\postextual

% Draft only: list every .bib entry so unchecked references stay visible.
% TODO: remover \nocite{*} antes da submissão
\nocite{*}
\printbibliography[title={Referências}]

\end{document}
